\chapter{Sincronização, Calibração e Localização}

White Rabbit ou IEEE 1588 HA pode distribuir tempo e frequência determinísticos por fibra, mas sincronização temporal não equivale automaticamente à coerência de fase da portadora em Ka. A incerteza de fase é:
\[
\sigma_\phi=2\pi f\sigma_t.
\]
Cada amostra deve ser associada a uma época comum e corrigida pelo atraso medido da cadeia:
\[
t_{i,j,n}=T_{\mathrm{epoch}}+\frac{n}{f_{s,i,j}}+\delta_{i,j},\qquad t_i^\ast=t_i-\tau_i(f,T).
\]
Para iluminador $T$, alvo $x$ e receptor $R_i$, o excesso de caminho biestático é:
\[
\Delta\rho_i=\lVert x-T\rVert+\lVert x-R_i\rVert-\lVert T-R_i\rVert=c\Delta\tau_i.
\]
AOA, TDOA e FDOA são fundidos com covariância, modelos de propagação e diluição geométrica de precisão.

\section{Referência temporal GNSS}
Cada estação inclui provisoriamente pelo menos duas antenas GNSS externas (\texttt{TIM-001}), com centros de fase levantados e visão desobstruída do céu, afastadas das Yagis VHF/UHF e de outras fontes RF intensas. Um receptor multiconstelação produz UTC, 1PPS, referência de frequência e efemérides. As saídas propostas de 1PPS e 10 MHz (\texttt{TIM-002}) disciplinam um relógio atômico ou oscilador de alta estabilidade local; o relógio distribui então o sincronismo determinístico por White Rabbit a todas as faces e ADCs.

O centro de fase da antena GNSS, o atraso do cabo, o atraso do receptor e o atraso da distribuição são medidos e versionados. O GNSS fornece época absoluta e disciplina de frequência de longo prazo, enquanto o holdover atômico local preserva a continuidade durante uma perda de GNSS. A calibração RF ponta a ponta continua obrigatória:
\[
t_i^\ast=t_i-\tau_{\mathrm{GNSS},i}-\tau_{\mathrm{cabo},i}-\tau_{\mathrm{RF},i}-\tau_{\mathrm{ADC},i}.
\]
As coordenadas conhecidas da estação permitem verificar a consistência entre posição levantada, geometria dos satélites e sincronismo medido da rede, mas não eliminam sozinhas multipercurso, erro do centro de fase ou atraso analógico.

\begin{figure}[H]
\centering
\includegraphics[width=0.95\textwidth]{fig07_sincronizacao.png}
\caption{Sincronização e orçamento de incerteza.}
\label{fig:fig07_sincronizacao}
\end{figure}

\begin{figure}[H]
\centering
\includegraphics[width=0.95\textwidth]{fig08_geometria_bistatica.png}
\caption{Geometria biestática.}
\label{fig:fig08_geometria_bistatica}
\end{figure}
